\setcounter{section}{4}

\Needspace{5\baselineskip}
\hypertarget{densenet}{%
\section{DenseNet}\label{densenet}}

在ResNet中，我们看到了跳跃连接（x + F(x)）的巨大威力。DenseNet的作者提出了一个更激进的想法：为了最大化网络中各层之间的信息流动，能否让每一层都直接连接到它后面的所有层？

在这种连接模式下，每一层都从所有 preceding layers（前面的层）中接收``集体知识''作为输入，并将其自身的特征图传递给所有 subsequent layers（后面的层）。

\Needspace{5\baselineskip}
\hypertarget{ux4e00ux6838ux5fc3ux7ec4ux4ef6ux5bc6ux96c6ux5757dense-blockux4e0eux8fc7ux6e21ux5c42transition-layer}{%
\subsection{一、核心组件：密集块（Dense Block）与过渡层（Transition Layer）}\label{ux4e00ux6838ux5fc3ux7ec4ux4ef6ux5bc6ux96c6ux5757dense-blockux4e0eux8fc7ux6e21ux5c42transition-layer}}

DenseNet的结构由两种核心模块交替组成。

\Needspace{4\baselineskip}
\begin{enumerate}
\def\labelenumi{\arabic{enumi}.}
\tightlist
\item
  密集块（Dense Block）
\end{enumerate}

密集块是DenseNet的核心构建单元。在一个密集块内部，每一层都直接与块内的所有后续层相连。

数学表达：设第~l~层的输出为~x\_l，则第~l~层的输入是前面所有层输出的拼接（Concatenation）：x\_l = H\_l({[}x\_0, x\_1, \ldots, x\_\{l-1\}{]})
其中~{[} \ldots{} {]}~表示在通道维度（Channel Dimension）~上进行拼接，H\_l(·)~代表一个复合函数（如BN-ReLU-Conv）。

\Needspace{5\baselineskip}
结构剖析：

复合函数~H\_l：~通常是一个标准的序列，包括批量规范化（BN）、ReLU激活函数~和~3x3卷积（Conv）。这个3x3卷积通常只产生很少的特征图（例如~k=32~个），这个数字~k~被称为增长率（Growth Rate）。

拼接操作：~每一层都会将其输出的~k~个特征图添加到整个块的``集体知识''中。因此，越靠后的层，其输入通道数会越多（因为它拼接了前面所有层的输出）。

局部连接：~密集连接只发生在密集块内部。不同密集块之间不直接连接。

\Needspace{4\baselineskip}
\begin{enumerate}
\def\labelenumi{\arabic{enumi}.}
\setcounter{enumi}{1}
\tightlist
\item
  过渡层（Transition Layer）
\end{enumerate}

由于密集块内的拼接操作会使特征图的通道数快速增长（例如，一个包含10层、增长率k=32的块，最后一层的输入通道数可能高达~3 + 10*32 = 323），为了控制模型的复杂度和计算量，需要在密集块之间插入过渡层。

作用：~压缩特征图的尺寸和通道数，进行降维。

\Needspace{5\baselineskip}
结构：~通常包含三个部分：

批量规范化（BN）

1x1卷积：~在DenseNet-C中引入，用于降低通道数。通常会将通道数减少一个比例（例如减少到原来的一半），这个比例是一个超参数。

2x2平均池化（Average Pooling）：~步幅为2，用于将特征图的高和宽减半，实现空间尺寸上的降采样。

\Needspace{5\baselineskip}
\hypertarget{ux4e8cux6574ux4f53ux7f51ux7edcux67b6ux6784}{%
\subsection{二、整体网络架构}\label{ux4e8cux6574ux4f53ux7f51ux7edcux67b6ux6784}}

一个完整的DenseNet的结构例如：7\emph{7步长为2的初始卷积------最大池化------密集层------过渡层（1}1卷积+2*2平均池化）------密集层------过渡层------密集层------过渡层------密集层------全局平均池化------分类器

\Needspace{5\baselineskip}
\hypertarget{ux4e09densenetux7684ux4f18ux52bfux4e0eux6311ux6218}{%
\subsection{三、DenseNet的优势与挑战}\label{ux4e09densenetux7684ux4f18ux52bfux4e0eux6311ux6218}}

\Needspace{5\baselineskip}
优势：

\begin{enumerate}
\def\labelenumi{\arabic{enumi}.}
\item
  强大的抗过拟合能力：~参数效率极高，参数量远少于同等性能的ResNet或VGG，因此在数据量不是特别巨大的情况下表现优异。
\item
  极大地改善了梯度流：~密集的连接使得梯度能够直接传播到早期层，网络非常容易训练。
\item
  鼓励特征重用：~每一层都可以直接利用网络中最底层的原始特征和所有中间层的复合特征，模型非常高效。
\item
  自带正则化效果：~其结构本身就有抑制过拟合的效果。
\end{enumerate}

\Needspace{5\baselineskip}
挑战与缺点：

\begin{enumerate}
\def\labelenumi{\arabic{enumi}.}
\item
  高内存消耗：~这是DenseNet最显著的缺点。由于需要保存所有中间特征图用于拼接，在训练过程中对GPU内存的消耗非常大。有许多研究致力于优化其内存效率。
\item
  潜在的计算开销：~大量的拼接操作也可能带来计算上的开销。
\item
  并非越深越好：~由于通道数的爆炸式增长，DenseNet的深度受到一定限制（通常最多几百层），不像ResNet可以轻松做到上千层。
\end{enumerate}

\Needspace{5\baselineskip}
\hypertarget{ux56dbux901aux9053ux6570ux7206ux70b8ux95eeux9898}{%
\subsection{四、通道数``爆炸''问题}\label{ux56dbux901aux9053ux6570ux7206ux70b8ux95eeux9898}}

\Needspace{5\baselineskip}
假设我们有一个DenseNet的密集块（Dense Block），并定义以下参数：

输入通道数：~假设进入第一个密集块的输入特征图有~C0 = 64~个通道。

增长率（Growth Rate）k：~这是DenseNet最重要的超参数之一。它定义了每个卷积层会产生多少新的特征图。假设~k = 32。这意味着每一层~H\_l（BN-ReLU-Conv）都会输出~32~个通道。

复合函数：~每一层都执行~H\_l:~BN → ReLU → 3×3 Conv~(输出~k~个通道)。

第0层（输入层）：输入到密集块的特征图：{[}通道数 = C0 = 64{]}

第1层：对这个~64~通道的输入进行~H\_1~操作（3x3卷积），并产生~k = 32~个新的特征图。

当前总特征图数：~64 (原有的) + 32 (新增的) = 96~个通道。

第2层：对这~96~个通道的输入进行~H\_2~操作，产生~32~个新的特征图。

当前总特征图数：~96 + 32 = 128~个通道。

第3层：~产生~32~个新特征图。

当前总特征图数：~128 + 32 = 160~个通道。

以此类推\ldots 第~l~层：输入通道数：~C0 + k × (l - 1)

对于一个包含~L~层的密集块，其最终输出的总通道数为：C\_\{output\} = C0 + k * L

虽然从数学上看是线性增长，但在深度学习实践中，这种增长带来的后果是``爆炸式''的

内存占用爆炸：若网络有~L~层，则第~l~层的输出特征图需被后续所有层复用，主流深度学习框架（如PyTorch、TensorFlow）在实现拼接操作（concat）时，会为每次拼接开辟独立的内存空间存储拼接结果，意味着它需在内存中保存~L−l+1~次。因此，全网络总内存消耗近似为O(L2)（平方级复杂度），而传统网络如ResNet仅为~O(L)（线性复杂度）。反向传播需依赖前向计算的中间特征图计算梯度。DenseNet因密集连接，需存储所有中间特征图，保存这些中间特征图用于反向传播会消耗巨大的GPU内存，成为训练DenseNet的主要瓶颈。

计算量增长：~尽管每一层只产生~k=32~个通道，但计算这些输出所需的卷积操作（Conv2D）的输入通道数却在不断变大，计算开销也随之增加。

\Needspace{5\baselineskip}
DenseNet的解决方案：

\begin{enumerate}
\def\labelenumi{\arabic{enumi}.}
\item
  过渡层（Transition Layer）。1x1 卷积：~在DenseNet-C中引入，用于压缩通道数。例如，如果上一个密集块输出~C~个通道，1x1卷积会将其减少到~θC~个通道，其中~θ~是压缩因子（通常~θ=0.5）。这直接解决了通道数过多的问题。2x2 平均池化：将特征图的高和宽减半，进一步降低计算量。
\item
  更多优化
\end{enumerate}

\Needspace{8\baselineskip}
（1）内存高效实现\hspace{0pt}

共享缓存机制：预分配一块固定缓存，供所有拼接层复用，避免重复开辟内存，将显存占用从~O(L2)~降至~O(L)。

梯度重构技术：反向传播时动态重建中间特征，减少存储需求（详见论文~Memory-Efficient Implementation of DenseNets）。

\Needspace{8\baselineskip}
（2）结构压缩技术\hspace{0pt}

Bottleneck层（DenseNet-B）：在3×3卷积前添加1×1卷积压缩通道，降低计算量。

Transition层压缩（DenseNet-C）：通过压缩因子~θ（如0.5）减少特征图通道数。

DenseNet-BC：结合Bottleneck与Transition压缩，综合优化内存与计算效率。

\Needspace{8\baselineskip}
（3）替代架构设计\hspace{0pt}

VoVNet（One-Shot Aggregation）：仅在模块最后一层聚合前面所有层特征，保留多尺度特征的同时，将MAC降至与ResNet相当。

MFM（Max-Feature-Map）并不是结构化通道剪枝。它把特征图分组后逐元素取最大值，通过竞争式激活保留特征，并可减少输出通道数。

\FloatBarrier
\Needspace{5\baselineskip}
\hypertarget{ux53c2ux8003ux6587ux732e-15}{%
\subsection{参考文献}\label{ux53c2ux8003ux6587ux732e-15}}

\vspace{0.25em}
\begin{itemize}
\tightlist
\item
  Huang, G., Liu, Z., van der Maaten, L., \& Weinberger, K. Q. (2017). \href{https://arxiv.org/abs/1608.06993}{Densely Connected Convolutional Networks}. CVPR 2017.
\item
  Pleiss, G., Chen, D., Huang, G., Li, T., van der Maaten, L., \& Weinberger, K. Q. (2017). \href{https://arxiv.org/abs/1707.06990}{Memory-Efficient Implementation of DenseNets}. arXiv:1707.06990.
\item
  Lee, Y., Hwang, J., Lee, S., Bae, Y., \& Park, J. (2019). \href{https://openaccess.thecvf.com/content_CVPRW_2019/html/CEFRL/Lee_An_Energy_and_GPU-Computation_Efficient_Backbone_Network_for_Real-Time_Object_CVPRW_2019_paper.html}{An Energy and GPU-Computation Efficient Backbone Network for Real-Time Object Detection}. CVPR Workshops 2019.（VoVNet）
\item
  Wu, X., He, R., Sun, Z., \& Tan, T. (2018). \href{https://arxiv.org/abs/1511.02683}{A Light CNN for Deep Face Representation with Noisy Labels}. \emph{IEEE Transactions on Information Forensics and Security}, 13(11), 2884--2896.（MFM）
\end{itemize}

\clearpage
